\section*{第 \ref{ch:g7:digestion-and-nutrients} 章——消化与营养}

\begin{solution}{exo:g7:digestion-and-nutrients:1}
它们以\emph{化学}方式分解食物：把其中的大块成分——
淀粉、构建物质和脂肪——切成
便于进入\omterm{prop:g4:journey-of-food:absorb}{血液}的小单位。牙齿咀嚼和搅拌只是让食物碎块变小；
\omterm{def:g4:journey-of-food:digestion}{消化液}则改变了这些碎块的本质。
\end{solution}

\begin{solution}{exo:g7:digestion-and-nutrients:2}
两支试管各盛淀粉糊，置于体温下：
一支加入新鲜\omterm{def:g4:journey-of-food:tract}{唾液}，另一支加入煮沸过的\omterm{def:g4:journey-of-food:tract}{唾液}。
一小时后，新鲜\omterm{def:g4:journey-of-food:tract}{唾液}那支试管中的淀粉已变为糖，
煮沸\omterm{def:g4:journey-of-food:tract}{唾液}那支则未变。由此得出结论：
\omterm{def:g4:journey-of-food:digestion}{消化液}本身即可分解淀粉，无需口腔参与；
煮沸则破坏其有效成分。
\end{solution}

\begin{solution}{exo:g7:digestion-and-nutrients:3}
糖类（\omterm{prop:g7:organs-at-work:bill}{葡萄糖}）——来自淀粉类和甜味食物，如面包；
构建物质单位——来自肉类、\omterm{prop:g3:sorting-living-things:groups}{鱼类}、蛋类、奶制品和豆类；
脂肪单位——来自油类和黄油；
此外还有维生素、\omterm{prop:g4:what-plants-need:needs}{无机盐}和水，
主要来自水果和蔬菜等食物。
\end{solution}

\begin{solution}{exo:g7:digestion-and-nutrients:4}
维生素、\omterm{prop:g4:what-plants-need:needs}{无机盐}和水原样穿过肠壁。
膳食纤维则完全不受\omterm{def:g4:journey-of-food:digestion}{消化液}作用：未经分解，
它无法穿过肠壁，
而是作为残渣中有益的粗渣继续前行。
\end{solution}

\begin{solution}{exo:g7:digestion-and-nutrients:5}
\omterm{def:g4:journey-of-food:tract}{小肠}内壁上的指状突起，
高约 1 毫米，数以百万计。每个突起都包着
自身细小的\omterm{def:g4:heart-and-blood:vessels}{血管}。
\end{solution}

\begin{solution}{exo:g7:digestion-and-nutrients:6}
送往\omterm{prop:g7:digestion-and-nutrients:absorption}{肝脏}——这是负责分类和储存的器官。其一项作用是
将这一餐中的部分\omterm{prop:g7:organs-at-work:bill}{葡萄糖}储存起来，
留待日后释放。
\end{solution}

\begin{solution}{exo:g7:digestion-and-nutrients:7}
运输方面：食物中的大块成分无法穿过肠壁，
好比衣柜塞不进信箱。
在重组方面：机体需要用小单位搭建自身物质，
因此需要的是零件，而不是成套家具——
既便于运输，也便于重新组装。
\end{solution}

\begin{solution}{exo:g7:digestion-and-nutrients:8}
这是公平实验中的对照：除\omterm{def:g4:journey-of-food:digestion}{消化液}的活性状态外，
其余\omterm{def:g6:exploring-our-environment:conditions}{条件}完全相同。它那未变的淀粉
证实了分解作用确实来自新鲜\omterm{def:g4:journey-of-food:tract}{唾液}，
同时说明其有效成分不耐热，
煮沸后即遭破坏，宛如精密机械。
\end{solution}

\begin{solution}{exo:g7:digestion-and-nutrients:9}
三者都是大面积、薄、湿润且布满\omterm{def:g4:heart-and-blood:vessels}{血管}的交换表面，
负责运输：\omterm{def:g7:what-is-respiration:respiration}{氧气}和\omterm{def:g7:what-is-respiration:respiration}{二氧化碳}供肺泡与鳃丝使用；
\omterm{def:g7:digestion-and-nutrients:families}{营养物质}则供\omterm{prop:g7:digestion-and-nutrients:absorption}{绒毛}使用。
三者结构原理相同，只是位置不同。
\end{solution}

\begin{solution}{exo:g7:digestion-and-nutrients:10}
卵属于构建物质类食物。分解始于
胃中的酸性\omterm{def:g4:journey-of-food:digestion}{消化液}（即玻璃肠模型中方块的遭遇），
最终在\omterm{def:g4:journey-of-food:tract}{小肠}中完成，生成小的构建物质单位。
穿越途径：\omterm{prop:g7:digestion-and-nutrients:absorption}{绒毛}、\omterm{prop:g4:journey-of-food:absorb}{血液}、\omterm{prop:g7:digestion-and-nutrients:absorption}{肝脏}。输送去向：
愈合中的伤口，分裂的\omterm{def:g1:the-five-senses:sense}{皮肤}\omterm{def:g5:first-look-at-cells:cell}{细胞}以这些单位
作为重建的材料。
\end{solution}

\begin{solution}{exo:g7:digestion-and-nutrients:11}
构建物质单位供应不足。首先受影响的是
一切处于建造与更新中的过程：生长本身、
\omterm{def:g3:bones-and-muscles:muscle}{肌肉}的维持、修复，以及不断更新的
内膜和\omterm{prop:g4:journey-of-food:absorb}{血液}成分——这些都无法由燃料类\omterm{def:g7:digestion-and-nutrients:families}{营养物质}替代，
因为砖块与柴火
并不能互换。
\end{solution}

\begin{solution}{exo:g7:digestion-and-nutrients:12}
举例来说：一餐食物以“成套家具”的形态进入\omterm{def:g4:journey-of-food:tract}{消化道}——淀粉、
构建物质和脂肪。沿途各段\omterm{def:g4:journey-of-food:digestion}{消化液}依次将其切分：
糖类、构建物质单位和脂肪单位，即送货清单。
在\omterm{prop:g7:digestion-and-nutrients:absorption}{绒毛}处，清单进入\omterm{prop:g4:journey-of-food:absorb}{血液}；\omterm{prop:g7:digestion-and-nutrients:absorption}{肝脏}负责分类、储存与释放，
\omterm{prop:g4:journey-of-food:absorb}{血液}随即分送：燃料送往工作中的\omterm{def:g3:bones-and-muscles:muscle}{肌肉}，
砖块送往建筑工地，其余送入储备。
盘中之食，
几小时内便化作带有地址的快递，
在体内循环流转。
\end{solution}

\begin{solution}{pb:g7:digestion-and-nutrients:1}
\textbf{1.} M：淀粉消失，糖类存在。M$'$：淀粉
未变。S：方块被\omterm{def:g4:journey-of-food:digestion}{消化}——缩小、变软。
S$'$：方块未变。

\textbf{2.} \omterm{def:g4:journey-of-food:tract}{唾液}中的有效成分把淀粉分解为
糖类。M$'$ 为对照：除\omterm{def:g4:journey-of-food:tract}{唾液}已煮沸外，
其余\omterm{def:g6:exploring-our-environment:conditions}{条件}与 M 完全相同，证明分解淀粉的
只能是新鲜\omterm{def:g4:journey-of-food:tract}{唾液}，试管中别无他物。

\textbf{3.} 它已开始分解卵中的构建物质，
将其切割为较小的单位，即构建物质类\omterm{def:g7:digestion-and-nutrients:families}{营养
物质}。（变透明、缩小的方块，正是这一切割
过程的可见表现。）

\textbf{4.} 真正起作用的是化学过程：各试管中既无牙齿咀嚼，
也无搅动，只要恰当的\omterm{def:g4:journey-of-food:digestion}{消化液}遇到恰当的食物，分解便会发生。
机械破碎
只是服务于化学过程。

\textbf{5.} \omterm{def:g4:journey-of-food:digestion}{消化液}是有工作\omterm{def:g6:exploring-our-environment:conditions}{条件}的活性物质，
\omterm{def:g6:exploring-our-environment:conditions}{条件}：遇冷便停止工作，——与从\omterm{def:g6:decomposers-and-soil:decomposers}{分解者}到\omterm{ex:g6:producing-our-food:bread}{酵母}的每一位工作者一样。
试管中的食物将不被\omterm{def:g4:journey-of-food:digestion}{消化}，
这只能说明冰箱会使化学反应暂停。

\textbf{6.} 每种\omterm{def:g4:journey-of-food:digestion}{消化液}各有其特性，以及各自所在部位的
工作\omterm{def:g6:exploring-our-environment:conditions}{条件}：若混于同一试管，二者都无法
发挥最佳效果，模型也就失去了\omterm{def:g4:journey-of-food:tract}{消化道}
逐段有序的顺序——而这正是模型
应当说明的内容。

\textbf{7.} 吸收壁：一个布满\omterm{prop:g7:digestion-and-nutrients:absorption}{绒毛}的
表面。其材料须薄，以便物质易于穿过，
且对小分子单位有选择性地通透——
在体内，其后有\omterm{def:g4:heart-and-blood:vessels}{血管}将穿过的物质带走，
将穿过者运走。

\textbf{8.} 未变化的薄片：我们所用的\omterm{def:g4:journey-of-food:digestion}{消化液}均不切割
纤维。其含义是：模型正确预测了纤维的
真实去向——不穿过任何结构，而是随之推进——
这一例外恰好标示出系统的局限。

\textbf{9.} 穿过\omterm{prop:g7:digestion-and-nutrients:absorption}{绒毛}的薄壁；进入
\omterm{prop:g4:journey-of-food:absorb}{血液}；经过\omterm{prop:g7:digestion-and-nutrients:absorption}{肝脏}，\omterm{prop:g7:digestion-and-nutrients:absorption}{肝脏}储存其中一部分；随\omterm{prop:g4:journey-of-food:absorb}{血液}
普遍输送至全身\omterm{def:g5:first-look-at-cells:cell}{细胞}——正在工作的\omterm{def:g3:bones-and-muscles:muscle}{肌肉}
在一小时内将其燃烧。

\textbf{10.} \omterm{def:g6:exploring-our-environment:conditions}{条件}是机制的组成部分：每个
部位都经过调节——此处为酸性，彼处被中和——因此
\omterm{def:g4:journey-of-food:digestion}{消化液}必须按正确的顺序与食物相遇，这正是
吞咽后食物依次经过各部位所提供的\omterm{def:g6:exploring-our-environment:conditions}{条件}。
\omterm{def:g4:journey-of-food:tract}{消化道}是一条生产线，而不是一只桶。

\textbf{11.} 因为\omterm{def:g4:journey-of-food:digestion}{消化}的工作者是物质，
而不是器官：如试管架所示，\omterm{def:g4:journey-of-food:digestion}{消化液}在恰当\omterm{def:g6:exploring-our-environment:conditions}{条件}下与食物相遇处，
便会分解食物。若将所缺的\omterm{def:g4:journey-of-food:digestion}{消化液}送至恰当的
部位，胶囊便可完成该部位的化学作用。
部位，胶囊便可完成该部位的化学作用。

\textbf{12.} “各\omterm{def:g4:journey-of-food:digestion}{消化液}在各自的部位将食物
切割为\omterm{def:g7:digestion-and-nutrients:families}{营养物质}——对照试管证明，切割食物的
正是\omterm{def:g4:journey-of-food:digestion}{消化液}，而别无他物。”
\end{solution}
